% ---
% title: The lognormal stock price
% tags: finance, function-of-rv, normal, distribution-function
% ---

\begin{problem}
	\begin{context}
		If a stock costs $S_0$ today and $S$ in a year, its \emph{log return}
		over the year is $X=\ln(S/S_0)$, so that $S=S_0\,e^{X}$. Log returns are
		convenient: the log return over two years is the sum of the two yearly
		log returns. The standard model of stock prices, which is also behind the
		Black--Scholes formula for option prices, assumes that the log return is
		normally distributed. Then the price $S=S_0e^X$ is said to have a
		\emph{lognormal} distribution: its logarithm is normal. Unlike a normal
		random variable, the price can never become negative.
	\end{context}
	A stock costs $S_0=100$ euros today. Its log return over the next year is
	$X\sim N(0.06,\,0.2^2)$, and its price in one year is $S=100\,e^{X}$.
	(a) Find the distribution function $F_S(s)=P(S\le s)$ of $S$, in terms of
	the standard normal distribution function $\Phi$.
	(b) Find the probability that the price in one year is below today's
	price, $P(S<100)$.
	(c) Find the median of $S$: the price $m$ with $P(S\le m)=0.5$.
	(d) Find the price level $s$ with $P(S\le s)=0.05$, using
	$z_{0.95}=1.645$. Interpret it.
\end{problem}

\begin{slideproblem}
	Price today $S_0=100$; log return $X\sim N(0.06,\,0.2^2)$; price in a year
	$S=100\,e^{X}$ (lognormal).
	(a) Find $F_S(s)=P(S\le s)$ with $\Phi$.
	\ (b) $P(S<100)$?
	\ (c) The median $m$ of $S$: $P(S\le m)=0.5$.
	\ (d) The level $s$ with $P(S\le s)=0.05$ \ ($z_{0.95}=1.645$).
\end{slideproblem}

\begin{slides}
	\small
	\exstatement
	\pause
	\sol (a) $S>0$, so $F_S(s)=0$ for $s\le0$. For $s>0$, $\ln$ is
	increasing:
	\[ F_S(s)=P\big(100e^X\le s\big)=P\Big(X\le\ln\frac{s}{100}\Big)
		=\Phi\Big(\frac{\ln(s/100)-0.06}{0.2}\Big). \]
	\pause
	(b) $P(S<100)=P(X<\ln1)=P(X<0)=\Phi\big(\frac{0-0.06}{0.2}\big)
		=\Phi(-0.3)=1-0.6179=\alert{0.3821}$.
\end{slides}

\begin{slides}
	\small
	\begin{columns}[T]
		\begin{column}{.55\textwidth}
			(c) $F_S(m)=0.5=\Phi(0)$ means
			$\ln(m/100)=0.06$:
			\[ m=100\,e^{0.06}\approx\alert{106.18}. \]
			\pause
			(d) $\Phi(z)=0.05$ for $z=-1.645$:
			\[ \ln\frac{s}{100}=0.06-1.645\cdot0.2=-0.269, \]
			\[ s=100\,e^{-0.269}\approx\alert{76.41}. \]
			\pause
			With 95\% probability the price stays above $76.41$: the loss on
			one share is at most $23.59$ euros.
		\end{column}
		\begin{column}{.43\textwidth}
			\centering
			\begin{tikzpicture}[x=.026cm,y=110cm]
				\fill[teal!30] (40,0) -- plot[domain=40:100,samples=50]
					(\x,{exp(-((ln(\x/100)-0.06)/0.2)^2/2)/(0.2*\x*2.5066)})
					-- (100,0) -- cycle;
				\fill[gold!60] (40,0) -- plot[domain=40:76.41,samples=30]
					(\x,{exp(-((ln(\x/100)-0.06)/0.2)^2/2)/(0.2*\x*2.5066)})
					-- (76.41,0) -- cycle;
				\draw[->] (40,0) -- (210,0) node[right] {\scriptsize $s$};
				\draw[navy,thick] plot[domain=40:200,samples=80]
					(\x,{exp(-((ln(\x/100)-0.06)/0.2)^2/2)/(0.2*\x*2.5066)});
				\draw[navy,dashed] (106.18,0) -- (106.18,.021)
					node[above] {\tiny median};
				\foreach \x in {76,100,150,200}
				\draw (\x,.0006) -- (\x,-.0006) node[below] {\tiny $\x$};
				\node[teal!60!black] at (160,.012) {\tiny $f_S(s)$};
				\node at (90,.0035) {\tiny (b)};
				\node at (58,.0032) {\tiny (d)};
			\end{tikzpicture}
		\end{column}
	\end{columns}
\end{slides}

\begin{solution}
	\idea $S$ is a function of the normal random variable $X$. Write every
	event about $S$ as an event about $X$, by taking logarithms; the logarithm
	is increasing, so the inequalities keep their direction. Then standardize
	$X$: if $X\sim N(\mu,\sigma^2)$, then $Z=(X-\mu)/\sigma$ is standard
	normal, so $P(X\le x)=\Phi\big((x-\mu)/\sigma\big)$, where $\Phi$ is the
	standard normal distribution function. Here $\mu=0.06$ and $\sigma=0.2$.

	\step{(a)} The price $S=100e^X$ is always positive, so $F_S(s)=0$ for
	$s\le0$. Take $s>0$. Divide by 100, then take the natural logarithm of
	both sides; as $\ln$ is increasing, the event does not change:
	\[ \{S\le s\}=\big\{100e^X\le s\big\}=\Big\{e^X\le\frac{s}{100}\Big\}
		=\Big\{X\le\ln\frac{s}{100}\Big\}. \]
	Standardizing $X$,
	\[ F_S(s)=P\Big(X\le\ln\frac{s}{100}\Big)
		=P\bigg(Z\le\frac{\ln(s/100)-0.06}{0.2}\bigg)
		=\Phi\bigg(\frac{\ln(s/100)-0.06}{0.2}\bigg),\qquad s>0. \]
	This is the distribution function of the lognormal distribution. (Its
	density, the derivative of $F_S$, is drawn in the picture below.)

	\step{(b)} A price below 100 means a negative log return: $S<100$ exactly
	when $X<\ln(100/100)=\ln1=0$. So
	\[ P(S<100)=P(X<0)=\Phi\Big(\frac{0-0.06}{0.2}\Big)=\Phi(-0.3). \]
	By the symmetry of the standard normal density, $\Phi(-z)=1-\Phi(z)$,
	and the table gives $\Phi(0.3)=0.6179$, so
	\[ P(S<100)=1-0.6179=0.3821. \]
	Even though the expected log return is positive, the price falls over
	the year with probability about 38\%.

	\step{(c)} We need $F_S(m)=0.5$. Since $\Phi(0)=0.5$ (half of the
	standard normal probability lies on each side of 0), the argument of
	$\Phi$ in (a) must be $0$:
	\[ \frac{\ln(m/100)-0.06}{0.2}=0\iff\ln\frac{m}{100}=0.06
		\iff m=100\,e^{0.06}\approx106.18. \]
	So the price is above 106.18 euros with probability $\frac12$ and below
	it with probability $\frac12$. The median of the log return is $0.06$,
	and the exponential function, being increasing, carries the median over.

	\step{(d)} We need $F_S(s)=0.05$, that is, $\Phi(z)=0.05$ for the
	standardized value $z$. We know $\Phi(1.645)=0.95$, so by symmetry
	$\Phi(-1.645)=1-0.95=0.05$, and $z=-1.645$. Hence
	\[ \frac{\ln(s/100)-0.06}{0.2}=-1.645\iff
		\ln\frac{s}{100}=0.06-1.645\cdot0.2=-0.269, \]
	and
	\[ s=100\,e^{-0.269}\approx76.41. \]
	Interpretation: with probability 95\% the price in one year is above
	76.41 euros, so the loss on one share is at most $100-76.41=23.59$ euros
	with probability 95\%. This is the 95\% Value at Risk of one share over
	one year, in the lognormal model.

	In the picture the area under the density $f_S$ to the left of 100 is
	$0.3821$ (b), and the area to the left of 76.41 is $0.05$ (d). The
	density is skewed to the right: large prices are possible, but the
	price cannot go below 0.
	\begin{center}
		\begin{tikzpicture}[x=.04cm,y=160cm]
			\fill[teal!30] (30,0) -- plot[domain=30:100,samples=60]
				(\x,{exp(-((ln(\x/100)-0.06)/0.2)^2/2)/(0.2*\x*2.5066)})
				-- (100,0) -- cycle;
			\fill[gold!60] (30,0) -- plot[domain=30:76.41,samples=40]
				(\x,{exp(-((ln(\x/100)-0.06)/0.2)^2/2)/(0.2*\x*2.5066)})
				-- (76.41,0) -- cycle;
			\draw[->] (30,0) -- (225,0) node[right] {\small $s$};
			\draw[navy,thick] plot[domain=30:215,samples=100]
				(\x,{exp(-((ln(\x/100)-0.06)/0.2)^2/2)/(0.2*\x*2.5066)});
			\draw[navy,dashed] (106.18,0) -- (106.18,.021)
				node[above] {\scriptsize median $106.18$};
			\foreach \x in {50,100,150,200}
			\draw (\x,.0004) -- (\x,-.0004) node[below] {\scriptsize $\x$};
			\draw (76.41,.0004) -- (76.41,-.0004) node[below] {\scriptsize $76.41$};
			\node[navy] at (170,.011) {\small $f_S(s)$};
			\node at (90,.004) {\scriptsize $0.3821$};
			\node (q) at (45,.007) {\scriptsize $0.05$};
			\draw[->] (q) -- (71,.0012);
		\end{tikzpicture}
	\end{center}
	\begin{answer}
		(a) $F_S(s)=\Phi\big((\ln(s/100)-0.06)/0.2\big)$ for $s>0$, and $0$ for
		$s\le0$
		\quad (b) $P(S<100)=\Phi(-0.3)=0.3821$
		\quad (c) $m=100e^{0.06}\approx106.18$
		\quad (d) $s=100e^{-0.269}\approx76.41$: with probability 95\% the price
		stays above 76.41 euros
	\end{answer}
	\begin{remark}
		The expected price is not the median. One can show that
		$\E[S]=S_0\,e^{\mu+\sigma^2/2}=100\,e^{0.08}\approx108.33$, larger than
		$m\approx106.18$: the long right tail of the lognormal density pulls
		the mean up.
	\end{remark}
\end{solution}
